\begin{frame}{Final Demonstration}

    \begin{center}

        \begin{tikzpicture}[
            font=\small,
            demo/.style={
                draw=WorkshopBlue,
                fill=WorkshopLightBlue,
                rounded corners=5pt,
                minimum width=2.7cm,
                minimum height=0.9cm,
                align=center
            }
        ]

            \node[demo] (add) {
                \textbf{Add}\\
                Student
            };

            \node[
                demo,
                right=0.45cm of add
            ] (view) {
                \textbf{View}\\
                Students
            };

            \node[
                demo,
                right=0.45cm of view
            ] (update) {
                \textbf{Update}\\
                Student
            };

            \node[
                demo,
                right=0.45cm of update
            ] (delete) {
                \textbf{Delete}\\
                Student
            };

            \node[
                draw=WorkshopBlue,
                fill=WorkshopGreen,
                rounded corners=6pt,
                minimum width=6.5cm,
                minimum height=0.9cm,
                below=0.8cm of view,
                font=\Large\bfseries
            ] (persist) {
                JSON Persistence
            };

            \draw[->, thick, WorkshopBlue]
                (add) -- (view);

            \draw[->, thick, WorkshopBlue]
                (view) -- (update);

            \draw[->, thick, WorkshopBlue]
                (update) -- (delete);

            \draw[->, thick, WorkshopBlue]
                (view.south) -- (persist.north);

        \end{tikzpicture}

    \end{center}

    \vfill

    \begin{center}
        \Large
        \textbf{From Python concepts to a working application}
    \end{center}

\end{frame}